% ---------------------------------------------------------------------------
\section{The text as a chain of glyphs: words, spaces and lines}

\subsection{A fingerprint that no language has}

The first fact, partly known, is that the Voynich text is too regular. The uncertainty about the next glyph (a classical
measure, called $h_2$) is 2.2 bits, against 2.6--3.3 in languages with an alphabet of similar size. Taking nine measures
together (predictability of glyphs, spaces, vocabulary, repetitions, similarity between words of the same page, links
between neighbouring words), the Voynich lies farther from the nearest language than the most isolated language of the
sample lies from its own neighbour (distance 9.9 against 4.1). The languages that ``most resemble'' it (Potawatomi,
Chinantec, Ewe, Malagasy) have nothing in common with one another: the Voynich has no natural linguistic family.

Other properties are normal: the number of different words (about 2{,}900 per 10{,}000 words, as in medieval technical
texts), Zipf's law, and the apparently weak ``syntax'' (which is the same as in Latin technical texts: it was the Bible,
a very formulaic text, that was the wrong benchmark).

\subsection{The space is a weak boundary inside a chain of glyphs}

This is the first new finding, and it is the key to many of the others.

\begin{itemize}
\item \textbf{The space can be guessed from the neighbouring glyphs.} Looking only at the two glyphs before and the two
  after a point, one predicts whether there is a space with an accuracy (F1) of 0.86, higher than in 97\% of languages
  (median 0.64). The spaces that transcribers mark as uncertain fall exactly where the rule is uncertain, and where ZL and
  Takahashi disagree.
\item \textbf{The space does not depend on the word.} In languages, when a space is ``optional'' (for instance in a
  compound), the choice depends strongly on the whole word. In the Voynich it does not: only the neighbouring glyphs and
  the length of the piece being written matter (the scribe breaks a long piece more often and avoids detaching pieces of
  1--3 glyphs).
\item \textbf{Different cuts yield real words.} Removing the spaces and putting them back at random with the rule of the
  neighbouring glyphs, half of the wrongly cut pieces are nevertheless words that exist in the Voynich (50\%), against
  6.5\% in languages (maximum 25\%).
\item \textbf{The vocabulary fills the possible forms.} Among the glyph sequences most probable under the rules of the
  system, 43\% are Voynich words; in languages 4\%. And the frequent words are exactly the most probable sequences
  (correlation 0.58, against 0.11 in languages).
\item \textbf{The rules between words mirror those inside words.} The preferences between the last glyph of a word and
  the first of the next resemble those between two consecutive glyphs inside a word (correlation 0.47), more than in all
  71 languages (at most 0.39, usually below 0.2). The same rule holds inside words written without a space: internal
  \eva{q}'s follow a glyph \eva{-y}, \eva{-o} or \eva{-d} in 28 cases out of 29.
\end{itemize}

\noindent In short, the text behaves as a \textbf{chain of glyphs} with its own sequence rules, in which spaces are
inserted where the chain allows it. Voynich ``words'' are not units chosen one by one as in a language; they are pieces
of a chain. This does not yet say what the chain represents.

\subsection{Junction and junction rules between neighbouring words}

\begin{itemize}
\item \textbf{A strong junction.} The last glyph of a word says a lot about the first glyph of the next: 0.18--0.19 bits
  beyond chance. Of 71 languages only six exceed the Voynich (among them Sanskrit, which writes the assimilations between
  words, and the Arabic of the Qur'an). In the gibberish handwritten by volunteers it is 0.014; in published generators
  between 0 and 0.08.
\item \textbf{Two precise rules.}
  \begin{itemize}
  \item \eva{qo-} or \eva{o-} before the gallows: \eva{qo-} after a word ending in \eva{-y}, \eva{-o}, \eva{-d}
    (59--64\% of the time), \eva{o-} after \eva{-n}, \eva{-r}, \eva{-s}, \eva{-m} (\eva{qo-} only 23--27\%).
  \item \eva{-l} or \eva{-r} at the end of a word: \eva{-r} before a word beginning with \eva{a-} (85\%), \eva{-l} before
    \eva{k-}, \eva{t-}, \eva{d-}, \eva{l-}, \eva{s-}, \eva{q-} (63--86\%).
  \item They have the shape of rules such as English \emph{a}/\emph{an} or French \emph{liaison}. They do not say which
    sounds lie behind them.
  \end{itemize}
\item \textbf{The scribe looks one word ahead.} When two words copied from the line above would jointly violate a rule,
  the scribe changes the piece easiest to change: for \eva{-l/-r} the end of the first word (59\% against 21\%), for
  \eva{qo-/o-} the start of the second (60\% against 18\%). So while finishing a word the scribe already knows the next
  one.
\item \textbf{And no further.} For a given word in between, the last glyph of a word says nothing about the first glyph of
  the word after next; in languages this link exists in 34 texts out of 42. The Voynich looks one word ahead and one
  back, no more.
\item \textbf{A rule of the system.} The rules hold equally in the three main hands, in all sections, in all bifolia, and
  with all three transliterations.
\end{itemize}

\subsection{The line is a closed unit}

\begin{itemize}
\item \textbf{The junction stops at the line break.} Between the last word of a line and the first of the line below the
  link is zero, against 0.19 inside the line. In the prose of languages, which breaks lines wherever it happens, the link
  carries over by half or entirely; it is close to zero only in texts in which each line is a verse.
\item \textbf{And at the gap of a drawing.} Where a line is interrupted by a plant, the link between the words on the two
  sides drops from 0.167 to 0.026. In a language the grammar would carry over the drawing. Words after the drawing begin
  like words at the start of a line.
\item \textbf{Line-edge forms.} For the same rest of the word, at the end of a line \eva{-m} rises by 15 percentage
  points and \eva{-r} falls by 12 (\eva{dar} in mid-line, \eva{dam} at line end); at the start of a line \eva{s-} and
  \eva{y-} rise by about 9--10 points in place of \eva{ch-} and \eva{k-}. In the last line of a paragraph, which does not
  reach the margin, \eva{-m} is rarer: it is a form tied to the margin.
\item \textbf{With no preceding word, \eva{o-} is written.} In the labels of drawings, which are isolated words,
  \eva{qo-} before the gallows is almost absent (0--5\% against 30--68\% in lines): \eva{qo-} arises only after a word
  written continuously. This was a prediction made in advance and confirmed on unseen data.
\end{itemize}
No published generator has both a chain with weak spaces and a closed line: generator U3 has the chain but lets it cross
the line break; the others have neither.

% ---------------------------------------------------------------------------
\section{The page and the book: copying from the line above, avoiding repetition}

\subsection{The scribe copies from the line immediately above}

\begin{itemize}
\item \textbf{Copying.} A word has an identical or nearly identical word in the two lines immediately above far more than
  chance ($z = 17.4$ for the same paragraph vocabulary), more than twice as much as in the generator of Timm and Schinner,
  which was built precisely on this idea.
\item \textbf{Only the line immediately above.} Copying comes from the line just above (strongly) and a little from the
  one before; from the third line up, nothing. In languages the ``recurrence'' of words is spread over 3--4 lines (it is
  the continuity of discourse, with a maximum at two lines); in the Voynich it is entirely at distance one. The Timm and
  Schinner generator draws from up to six lines above: this is the first clear difference with their theory.
\item \textbf{Pieces of line.} Two neighbouring copied words come, more than chance, from two neighbouring words of the
  line above, usually in the same order but also inverted.
\item \textbf{Not by columns.} The copied word does not sit in the same column as its source (the apparent alignment came
  from the line edges).
\item \textbf{It stops at the paragraph and the page.} The first line of a paragraph does not copy the last line of the
  previous paragraph, nor the first line of a page the last line of the previous page. Inside the paragraph, from the
  third line on, copying is constant.
\item \textbf{The copy adapts to the junction rules.} A word copied from the line above takes the form (\eva{qo-} or
  \eva{o-}, \eva{-l} or \eva{-r}) that agrees with its \emph{new} neighbour more than that of its source (effect $+0.49$
  against $+0.16$). The junction rules are applied at the moment of writing.
\item \textbf{Gibberish does not do it.} The meaningless text handwritten by volunteers does not copy from the lines above
  ($-0.001$); languages do, but with recurrence spread over several lines.
\end{itemize}

\subsection{Each line avoids starting like the line above}

\begin{itemize}
\item Shuffling the order of lines within paragraphs shows that the scribe \textbf{avoids} beginning a line with the same
  start as the previous line: \eva{qo-} ($z = -9.7$), \eva{o-} ($-7.0$), \eva{d-} ($-3.8$), \eva{ch-} ($-3.6$), \eva{y-}
  ($-3.5$). If the line above starts with \eva{qo-}, the line below takes \eva{o-} (the probability of \eva{qo-} drops by
  52 points).
\item It concerns the start of the line, not any visual edge: where lines resume to the right of a drawing, one under the
  other, the effect is absent. The line two above does not matter, nor the end of the line above, nor the right margin.
\item Ratio of observed to expected identical starts: Voynich 0.51; languages median 0.95 (only 3 of 57 below the
  Voynich); handwritten gibberish 1.32; generators 0.93--2.05.
\end{itemize}

\subsection{Paragraphs, pages, bifolia}

\begin{itemize}
\item \textbf{The paragraph opens with a large sign.} 83\% of paragraph-initial words begin with a gallows (\eva{k},
  \eva{t}, \eva{p}, \eva{f} or compounds), against 9\% elsewhere; among them \eva{p} dominates (54\%), which is rare in the
  rest of the text. The opening gallows is three times less tied to its word than usual: it is largely a mark of position.
\item \textbf{First and last lines have their own register.} The first line of a paragraph has longer words, more
  \eva{p}, \eva{f}, \eva{sh}, more final \eva{-y}; the last has fewer \eva{qo-} and fewer gallows, more \eva{ch-},
  \eva{s-}, \eva{y-}.
\item \textbf{\eva{-ey} increases down the page} ($+0.12$ to $+0.14$ from top to bottom, for the same position in the
  line), especially in the recipe pages. Lines get shorter down the page.
\item \textbf{Each page has a strong identity}, mainly along an axis from words in \eva{-edy/-eey} to words in
  \eva{-aiin/-ain} (the page departs from the profile of its section 2.8 times more than chance).
\item \textbf{The bifolium is a working unit.} The identity of a page is shared by the two sides of the same leaf and by
  the conjugate leaf of the same bifolium (which in the bound book lies far away), not by the facing page. The two halves
  of a bifolium share rare words and copy from each other. All 48 bifolia with enough data are entirely in language A or
  entirely in language B, and 94\% are in a single hand.
\item \textbf{Languages A and B differ in vocabulary}, not in glyphs: the best substitution of glyphs brings A closer to B
  by only 1\% of the gap, while the same search recovers a random key 100\% of the time. B is not A enciphered with
  another key.
\item \textbf{The order of the pages largely preserves the order of writing:} consecutive pages share more words than any
  two pages of the same section.
\end{itemize}

% ---------------------------------------------------------------------------
\section{Choices between variants: repeating, persisting, drifting}

In the Voynich many words exist in pairs of variants differing by one glyph: \eva{qo/o} at the start of a word,
\eva{k/t} (two gallows), \eva{sh/ch}, \eva{-ey/-dy} at the end. Studying \textbf{how the scribe chooses} between these
variants turned out to be the most fruitful way of looking at the text.

\subsection{Repeating immediately}

\begin{itemize}
\item In the Voynich 1.0\% of adjacent word pairs are the same word repeated (\eva{chol chol}, \eva{qokeedy qokeedy}) and
  3.9\% are a nearly identical word (one glyph different). In languages: 0.10\% and 0.23\% (medians; maxima 0.71\% and
  2.2\%); in the six Nordic scribes 0.01--0.05\% and 0.2--0.5\%; in the 73 German scribes median 0.05\% and 0.46\% (maxima
  0.54\% and 1.20\%); in handwritten medieval technical texts (herbals, recipe books, astrology) 0.03--0.15\% and
  0.2--0.6\%; in the Copiale copyist practically never.
\item Repetitions are the same in all sections (herbal 1.0\%, biological 1.2\%, pharmaceutical 0.9\%, recipes 0.9\%):
  they do not depend on the genre of the content.
\item \textbf{Repetition fades quickly:} within the line it is strong between nearby words and disappears after 6--7
  words (far/near ratio 0.12). In languages it is the opposite (median 2.86): languages avoid repeating at once, the
  Voynich repeats at once. No comparison text has this signature.
\item When a repeated word changes its start or end, the change follows the junction rule with its neighbour
  (\eva{okeey qokeey}: after \eva{-y} a \eva{q} is added).
\end{itemize}

\subsection{The ``state'' of the choices}

This is the finding that required most work, and its description was refined several times along the way (Appendix~B).
This is the final state.

\begin{itemize}
\item \textbf{The short state.} Between nearby words of the same line the scribe tends to make the same choice. Removing
  the preferences of each word, of the page and of the position, and correcting the bias of the measure, the agreement is
  $+0.13$ between adjacent words and $+0.09$ at two and three words (\eva{k/t}, \eva{sh/ch}, \eva{-ey/-dy} together;
  $+0.14$ for \eva{k/t} and $+0.18$ for \eva{-ey/-dy} alone). It halves in about three words.
\item \textbf{It is not copying.} Agreement is almost the same between very different words (four or more glyphs apart:
  $+0.09$) and between similar words ($+0.13$). The Timm and Schinner generator, which copies nearby words with
  modifications, does the opposite: agreement only between similar words.
\item \textbf{It is not position in the line.} With a comparison that preserves position (and a ``placebo'' to separate the
  effect of the procedure), the part due to position is about zero.
\item \textbf{It is separate for each choice.} Between different choices the agreement is zero ($-0.01$): there is no
  common ``marked mode''; \eva{k/t} and \eva{-ey/-dy} change each on its own.
\item \textbf{It does not depend on how many times the choice occurs in between:} it lives along the line, from word to
  word, and is ``read'' when the choice occurs.
\item \textbf{It is the same in all hands} (hands 1, 2, 3: between $+0.12$ and $+0.15$ between adjacent words), in
  languages A and B, in all sections where it can be measured, and in the three transliterations (in Glen Claston the
  gallows $+0.12$, the endings $+0.14$). It is a \textbf{rule of the writing system}, not the habit of one scribe.
\item \textbf{Variants of form alone show no clear agreement.} The two forms of \eva{d} and of \eva{e} that only Glen
  Claston distinguishes (differences in the shape of the sign, not of the word) have a weak and uncertain agreement
  ($+0.05$ / $+0.06$, intervals touching zero). We do not draw a general rule from this: a real scribe has a window in a
  choice of form alone (Section~7).
\item \textbf{A slow state, in addition.} Besides the short state there is a slow component that crosses the line break:
  words of consecutive lines agree ($+0.05$), at two lines less ($+0.03$), at 7--12 lines not at all. It is not a
  top-to-bottom trend of the page. Unlike the short component, this slow component \textbf{is normal in real scribes}
  (Section~7).
\end{itemize}

\paragraph{What it says and what it does not}
A state lasting a few words, separate for each choice, is compatible with a regulated writing habit (the writer ``enters''
a variant and stays in it for a while), with a rule of an encryption system, or with a procedure that produces text. It
does not say that the variants carry information, nor which.

\subsection{Drift along the line}

\begin{itemize}
\item Inside the line, for the same word and without the first and last word, the variants \eva{qo}, \eva{k}, \eva{sh},
  \eva{-ey} become rarer towards the right: from $-0.020$ to $-0.034$ per 10 glyphs, that is 8--12 percentage points from
  one end of a typical line to the other. The drift is steeper in the first half of the line and is present in every
  hand, about three times stronger in language B.
\item Being common to all choices, it is not an effect of ink on a single stroke (for example the small stroke that
  distinguishes \eva{sh} from \eva{ch}). No published generator has it.
\item Next to drawings the marked forms become rarer; in particular the \eva{q} drops in the word that begins right after
  a drawing ($-33$ points), not in the one before. Without the images we cannot tell whether this is writing (restarting
  against the drawing with \eva{o-}) or reading by the transcribers.
\end{itemize}
